\documentclass[11pt,a4paper]{article}
\usepackage[margin=2.3cm]{geometry}
\usepackage{amsmath,amssymb}
\usepackage[T1]{fontenc}
\usepackage[utf8]{inputenc}
\usepackage{lmodern}

\title{A $\tau$-First Symbolic Van Vleck Derivation of Quartic Watson Constants}
\author{Codex + custom SymPy workflow}
\date{\today}

\begin{document}
\maketitle

\section*{Summary}
We implemented a symbolic Van Vleck derivation of quartic centrifugal distortion constants
starting from the standard rovibrational block decomposition
\begin{equation}
H = H_{20} + H_{02} + H_{12} + H_{30} + H_{22} + H_{40},
\end{equation}
with
\begin{align}
H^{(0)} &= H_{20}+H_{02},\\
H^{(1)} &= H_{12}+H_{30},\\
H^{(2)} &= H_{22}+H_{40}.
\end{align}

The key point is that the derivation is stopped first at the symmetric quartic tensor
\begin{equation}
\tau=(\tau_{xxxx},\tau_{yyyy},\tau_{zzzz},\tau_{xxyy},\tau_{xxzz},\tau_{yyzz}),
\end{equation}
rather than directly at Watson constants. Watson A- and S-reduction constants are then obtained
by linear projection of \(\tau\).

\section*{Symbolic Filtering}
To keep the algebra tractable, each operator monomial is filtered according to:
\begin{itemize}
\item rotational degree \(\le 4\),
\item even bosonic degree,
\item allowed contribution class only:
\[
(H_{12},H_{12}),\ (H_{22}),\ (H_{12},H_{30}),\ (H_{30},H_{30}),\ (H_{40}).
\]
\end{itemize}

\section*{Second-Order Result}
At second order, only the class \((H_{12},H_{12})\) contributes to the quartic tensor. In
general tensor form, this means that \(\tau\) is bilinear in the first rovibrational coupling
tensor \(\mu_{\alpha\beta,k}\mu_{\gamma\delta,k}\), with the expected harmonic denominators.
In the monomial basis
\[
J_x^4,\ J_y^4,\ J_z^4,\ J_x^2J_y^2,\ J_x^2J_z^2,\ J_y^2J_z^2,
\]
the exact second-order coefficients are
\begin{align}
\tau_{\alpha\alpha\alpha\alpha}^{(2)} &=
-\sum_k \frac{\mu_{\alpha\alpha,k}^2}{8\omega_k^2},\\
\tau_{\alpha\alpha\beta\beta}^{(2)} &=
-\sum_k \frac{\mu_{\alpha\alpha,k}\mu_{\beta\beta,k}}{4\omega_k^2},
\qquad \alpha \neq \beta.
\end{align}
This is the standard known second-order result, and in the present implementation it serves
as the primary internal validation benchmark.

For the one-mode diagonal symbolic run used for explicit verification, the quartic tensor
coefficients are
\begin{align}
\tau_{xxxx} &= -\frac{\mu_{xx,1}^2}{8\omega_1^2}, &
\tau_{yyyy} &= -\frac{\mu_{yy,1}^2}{8\omega_1^2}, &
\tau_{zzzz} &= -\frac{\mu_{zz,1}^2}{8\omega_1^2},\\
\tau_{xxyy} &= -\frac{\mu_{xx,1}\mu_{yy,1}}{4\omega_1^2}, &
\tau_{xxzz} &= -\frac{\mu_{xx,1}\mu_{zz,1}}{4\omega_1^2}, &
\tau_{yyzz} &= -\frac{\mu_{yy,1}\mu_{zz,1}}{4\omega_1^2}.
\end{align}

Projection onto Watson A reduction gives
\begin{align}
D_J &= \frac{\mu_{xx,1}\mu_{yy,1}}{8\omega_1^2},\\
D_{JK} &= \frac{-2\mu_{xx,1}\mu_{yy,1} + \mu_{xx,1}\mu_{zz,1} + \mu_{yy,1}\mu_{zz,1}}{8\omega_1^2},\\
D_K &= \frac{\mu_{xx,1}\mu_{yy,1} - \mu_{xx,1}\mu_{zz,1} - \mu_{yy,1}\mu_{zz,1} + \mu_{zz,1}^2}{8\omega_1^2},\\
d_1 &= \frac{\mu_{zz,1}(\mu_{xx,1}-\mu_{yy,1})}{8\omega_1^2},\\
d_2 &= \frac{\mu_{xx,1}^2 - \mu_{yy,1}^2 - 2\mu_{zz,1}(\mu_{xx,1}-\mu_{yy,1})}{16\omega_1^2}.
\end{align}

Thus, at second order, the quartic Watson constants originate exclusively from the
\((H_{12},H_{12})\) channel.

\section*{Higher-Order Structure}
No new quartic ground-state contribution was found at third order in the current reduced run.
At fourth order, the expected new contributions appear from
\begin{equation}
(H_{22}),\qquad (H_{12},H_{30}),\qquad (H_{30},H_{30}),\qquad (H_{40}),
\end{equation}
which contribute first to the quartic tensor \(\tau\) and only afterwards project onto the
Watson quartic constants. In the present one-mode symbolic run, however, the explicit output
separates only into the two nonzero channels
\[
\tau^{(4)}=\tau[H_{30},H_{30}] + \tau[H_{12},H_{30}],
\]
with diagonal elements
\[
\tau_{\alpha\alpha\alpha\alpha}^{(4)}
=
-\frac{\hbar\,\mu_{\alpha\alpha}^{(1)}\mu_{\alpha\alpha}^{(1)}\left(\Phi^{(3)}\right)^2}{1728\,\omega^7}
+\frac{\hbar\,\mu_{\alpha\alpha}^{(1)}\mu_{\alpha\alpha}^{(2)}\Phi^{(3)}}{144\,\omega^5}.
\]

\section*{Conclusion}
The natural physical output of the symbolic Van Vleck derivation is the quartic tensor
\(\tau\). This separates the derivation itself from the later choice of Watson reduction,
representation, or projection convention.

\end{document}
